\documentclass[11pt]{article}

% ---------- Packages ----------
\usepackage[a4paper,margin=2.5cm]{geometry}
\usepackage[utf8]{inputenc}
\usepackage[T1]{fontenc}
\usepackage{amsmath}
\usepackage{graphicx}
\usepackage{xcolor}
\usepackage{tabularx}
\usepackage{booktabs}
\usepackage{enumitem}
\usepackage{hyperref}
\usepackage[explicit]{titlesec}
\usepackage{fancyhdr}
\usepackage{caption}
\usepackage{listings}

% ---------- Colors ----------
\definecolor{etsitorange}{RGB}{237,125,49}
\definecolor{etsitblue}{RGB}{31,73,125}
\definecolor{codebg}{RGB}{245,245,245}

% ---------- Code listing style ----------
\lstset{
  backgroundcolor=\color{codebg},
  basicstyle=\ttfamily\small,
  breaklines=true,
  frame=single,
  columns=fullflexible,
  keywordstyle=\color{etsitblue}\bfseries,
  showstringspaces=false
}

% ---------- Section banners (orange bar, like the template) ----------
\titleformat{\section}
  {\bfseries\Large\color{white}}
  {}{0pt}
  {\colorbox{etsitorange}{\parbox{\dimexpr\textwidth-4\fboxsep-12pt\relax}{\centering\thesection.\ #1}}}
\titlespacing*{\section}{0pt}{18pt}{12pt}

% Unnumbered \section* (used internally by \tableofcontents, \listoftables,
% \listoffigures) gets the same orange banner, without a number.
\titleformat{name=\section,numberless}
  {\bfseries\Large\color{white}}
  {}{0pt}
  {\colorbox{etsitorange}{\parbox{\dimexpr\textwidth-4\fboxsep-12pt\relax}{\centering#1}}}
\titlespacing*{name=\section,numberless}{0pt}{18pt}{12pt}

\titleformat{\subsection}
  {\bfseries\large}
  {\thesubsection.}{0.5em}{#1}
\titlespacing*{\subsection}{0pt}{12pt}{6pt}

\titleformat{\subsubsection}
  {\bfseries\normalsize}
  {\thesubsubsection}{0.5em}{#1}
\titlespacing*{\subsubsection}{0pt}{10pt}{4pt}

% Appendix banner (no numbering, e.g. "APPENDIX: ...")
\newcommand{\appendixbanner}[1]{%
  \colorbox{etsitorange}{\parbox{\dimexpr\textwidth-4\fboxsep-12pt\relax}{\centering\bfseries\Large\color{white} #1}}%
  \vspace{12pt}
}

\renewcommand{\contentsname}{CONTENTS}
\renewcommand{\listtablename}{LIST OF TABLES}
\renewcommand{\listfigurename}{LIST OF FIGURES}

\pagestyle{fancy}
\fancyhf{}
\renewcommand{\headrulewidth}{0pt}
\fancyfoot[C]{\thepage}

\usepackage{float}
\usepackage{array}
\usepackage{cite}
\hypersetup{hidelinks}
\graphicspath{{figures/}}
\setlength{\emergencystretch}{2em}  % avoid long \texttt paths running into the margin

% Numbers quoted in the text, generated from results/logs/*.json by report/make_numbers.py
\input{numbers}

% Claims still to be checked against a quoted source (redefine as empty for the final version)
\newcommand{\todoverify}[1]{\textcolor{red}{\,\textsf{\footnotesize[verify: #1]}}}

\begin{document}

% =====================================================================
% COVER PAGE
% =====================================================================
\begin{titlepage}
\centering
\vspace*{1cm}
\includegraphics[width=0.4\textwidth]{figures/etsit_logo.jpeg}\par
\vspace{1.5cm}

{\color{etsitblue}\bfseries\Large Course:}\par
\vspace{0.5cm}
{\bfseries\huge Machine Learning Lab}\par
\vspace{1cm}
{\Large Half-term Activities Report}\par
{\bfseries\Large Developing Machine Learning models}\par
\vspace{0.5cm}
\rule{0.7\textwidth}{0.4pt}\par
\vspace{3cm}

\large Name: \underline{\hspace{2mm}Oriana Andrea Ferraz De Jesus\hspace{2mm}}\par
\vspace{0.5cm}
\hfill Date: \underline{\hspace{2mm}October 2026\hspace{2mm}}\par

\end{titlepage}

% =====================================================================
% TABLE OF CONTENTS / LIST OF TABLES / LIST OF FIGURES
% =====================================================================
\tableofcontents
\vspace{1cm}

\listoftables
\vspace{1cm}

\listoffigures

\newpage

% =====================================================================
% 1. METHODOLOGY TO DEVELOP MACHINE LEARNING MODELS
% =====================================================================
\section{Methodology to develop Machine Learning models}

\subsection{Methodological guidelines for developing Machine Learning models}

The project predicts the grade of TCGA gliomas (LGG vs GBM) and the age at diagnosis from clinical
data and the mutation status of 20 genes, and explores the mutation profiles without labels (the
theory and the critical analysis are in the companion PRDL report). The development followed five
guidelines, each enforced by the code rather than by convention:

\begin{enumerate}
  \item \textbf{Staged, reproducible workflow.} Five notebooks run in order (data preparation, EDA,
        unsupervised analysis, regression, classification). Each one reads only the output of the
        previous stages, fixes the seed (42) for every random process, and writes its results to
        \texttt{results/}: figures, \LaTeX{} tables and a JSON summary per stage. Both reports are built
        from these files: every number in the text is a macro generated from the JSON summaries by
        \texttt{report/\allowbreak{}make\_\allowbreak{}numbers.\allowbreak{}py}, so no number is typed by hand.
  \item \textbf{Leakage prevention by construction.} Columns that encode the target (TCGA project,
        histological diagnosis) are excluded from the predictor lists in \texttt{src/data.py}, and the
        classification notebook stops if one of them appears. All preprocessing that learns from the
        data (standardisation, one-hot encoding, discretisation) is a step of a scikit-learn
        \texttt{Pipeline}, so it is refitted on every training fold \cite{kaufman2012}.
  \item \textbf{Nested cross-validation for every supervised model} (scheme of the course notebook
        \texttt{Example\_\allowbreak{}ExtrOSA\_\allowbreak{}ModelTuning\_\allowbreak{}NestedLoop}): 10 outer folds for evaluation and 3 inner
        folds for \texttt{GridSearchCV}, stratified for classification
        (Figure~\ref{fig:mllb-nested}).
  \item \textbf{Evaluation against a baseline, with uncertainty.} Out-of-fold predictions for every
        patient; bootstrap 95\% intervals; no model is declared better than another when the intervals
        overlap, and a paired bootstrap of the difference complements the comparison.
  \item \textbf{Verification of AI-generated code.} Every notebook was executed in a clean copy of
        the repository before delivery, every figure was inspected, and the run on the author's
        computer was compared value by value with that run (Appendix A).
\end{enumerate}

\begin{figure}[H]
\centering
\begin{tabular}{|p{0.92\textwidth}|}
\hline
\textbf{for} each outer fold $k = 1,\dots,10$ (stratified, seed 42):\\
\quad split patients into training part $T_k$ (90\%) and test fold $V_k$ (10\%)\\
\quad \textbf{for} each hyperparameter configuration $\theta$ in the grid:\\
\qquad 3-fold CV inside $T_k$: fit pipeline (preprocessing + model) on 2/3, score on 1/3\\
\quad choose $\theta^*_k$ with the best mean inner score (AUC or RMSE)\\
\quad refit the pipeline with $\theta^*_k$ on all of $T_k$\\
\quad predict $V_k$ $\rightarrow$ out-of-fold predictions; permutation importance on $V_k$\\
metrics and bootstrap intervals on the \NClean{} out-of-fold predictions\\
final model for interpretation: the same grid search applied once to all patients\\
\hline
\end{tabular}
\caption{Nested cross-validation as implemented in \texttt{src/\allowbreak{}evaluation.\allowbreak{}py} (\texttt{nested\_cv},
\texttt{fit\_final}).}
\label{fig:mllb-nested}
\end{figure}

\subsection{Tools to develop Machine Learning models}

Table~\ref{tab:mllb-tools} lists the software. The environment is created with Anaconda from
\texttt{environment.yml} (Python 3.11); all libraries are pinned in \texttt{requirements.\allowbreak{}txt} with
exact versions that were installed and tested together; the script \texttt{check\_env.py} verifies them
and writes the result to the
\texttt{results/logs/} folder. Notebooks were run in Visual Studio Code with a
dedicated Jupyter kernel; the code is versioned with git and published on GitHub
(Section~\ref{sec:mllb-code}).

\begin{table}[H]
\centering
\caption{Software used. Versions are those pinned in \texttt{requirements.\allowbreak{}txt}.}
\label{tab:mllb-tools}
\begin{tabularx}{\textwidth}{l l X}
\toprule
\textbf{Tool} & \textbf{Version} & \textbf{Use} \\
\midrule
Python & 3.11 & Language \\
pandas \cite{mckinney2010} & 2.2.3 & Data loading, cleaning and tables \\
NumPy \cite{harris2020} & 2.2.6 & Numerical arrays, bootstrap \\
SciPy \cite{virtanen2020} & 1.15.3 & Statistical tests, hierarchical clustering \\
scikit-learn \cite{pedregosa2011} & 1.6.1 & Pipelines, models, cross-validation, PCA, MDS, t-SNE, $k$-means \\
XGBoost \cite{chen2016} & 2.1.4 & Gradient boosting \\
statsmodels \cite{seabold2010} & 0.14.4 & OLS inference, VIF, influence measures, multiple testing \\
SHAP \cite{lundberg2017} & 0.48.0 & SHAP values of the winning models (TreeSHAP, linear) \\
Matplotlib \cite{hunter2007} / seaborn \cite{waskom2021} & 3.10.3 / 0.13.2 & Figures \\
Jupyter (ipykernel, nbconvert) & 6.29.5 / 7.16.6 & Notebooks and their execution \\
Anaconda, VS Code, git/GitHub & -- & Environment, editor, version control \\
Claude (Anthropic) & -- & AI code assistant (Appendix A) \\
\bottomrule
\end{tabularx}
\end{table}

The code is organised as notebooks that call functions from a small package, so that the same code
is used everywhere and can be read on its own during the defense:
\texttt{config.py} (paths, seed, validation settings), \texttt{data.py} (cleaning, encoding, feature
sets), \texttt{eda.py} (tests and effect sizes), \texttt{unsupervised.py}, \texttt{pipelines.py}
(preprocessing), \texttt{models.py} (pipelines and grids), \texttt{evaluation.py} (nested CV, metrics,
bootstrap, comparisons), \texttt{explain.py} (SHAP values and plots), \texttt{causal.py} (mediation
decomposition, robust OLS, causal diagram), \texttt{plotting.py} (style and fixed colours: LGG blue, GBM orange in every
figure) and \texttt{reporting.py} (\LaTeX{} tables and JSON summaries).

% =====================================================================
% 2. DATA PREPARATION
% =====================================================================
\section{Data Preparation}

\textbf{Input files} (\texttt{data/raw/}, never modified): the original TCGA export
\texttt{TCGA\_GBM\_LGG\_Mutations\_all.csv} (\NRaw{} patients, text labels, with patient identifier,
TCGA project and diagnosis) and the encoded UCI file \texttt{TCGA\_InfoWithGrade.csv} (\NClean{}
patients) \cite{tasci2022dataset}. Notebook \texttt{01\_data\_preparation} rebuilds the clean dataset
from the original export so that the identifier and the diagnosis are kept, in five steps:

\begin{enumerate}
  \item \textbf{Quality checks.} Missing values are not empty cells but the strings \texttt{--} and
        \texttt{not reported}, so \texttt{isna()} finds none; they are counted
        explicitly (function \texttt{count\_invalid\_markers} in \texttt{data.py}). No patient identifier is duplicated, and the gene columns contain only
        \texttt{MUTATED}/\texttt{NOT\_MUTATED}.
  \item \textbf{Exclusions} (Table~\ref{tab:mllb-cleaning}): each patient is assigned to the first
        reason that applies, so the counts add up; the excluded patients are \NRemovedLGG{} LGG and
        \NRemovedGBM{} GBM.
  \item \textbf{Encoding} with the UCI codes. Age is converted from ``\textit{X} years \textit{Y}
        days'' to $X + Y/365$ years, rounded to two decimals; race is kept as a nominal code and
        one-hot encoded later, inside each model pipeline.
  \item \textbf{Leakage columns.} \texttt{Project} is identical to the grade for all \NRaw{} patients
        and is dropped; \texttt{Primary\_Diagnosis} is kept only as metadata.
  \item \textbf{Check against UCI.} \texttt{data.verify\_against\_uci} compares the rebuilt dataset with
        the UCI file column by column and row by row (number of patients, order and every value of the
        24 columns); the notebook stops with an \texttt{assert} if any check fails. All checks pass.
\end{enumerate}

\begin{table}[H]
\centering
\caption{Patients removed at each cleaning step.}
\label{tab:mllb-cleaning}
\input{tables/data_cleaning.tex}
\end{table}

The resulting variables and their role in each task are listed in Table~\ref{tab:mllb-dictionary};
the predictor sets are defined once in \texttt{data.feature\_columns} and imported by every notebook
(23 predictors for classification, 22 for regression, the 20 genes for the unsupervised analysis).

\begin{table}[H]
\centering
\caption{Data dictionary of the clean dataset.}
\label{tab:mllb-dictionary}
\resizebox{\textwidth}{!}{\input{tables/data_dictionary.tex}}
\end{table}

% =====================================================================
% 3. EXPLORATORY DATA ANALYSIS (EDA)
% =====================================================================
\section{Exploratory Data Analysis (EDA)}

Notebook \texttt{02\_eda} follows the sections of the course EDA notebook (quality, target and
leakage, univariate analysis, normality, tests and effect sizes, outliers, correlation and VIF,
relevance, conditional probability, multivariate view, baseline) and adds an EDA of the regression
target. Notebook \texttt{03\_unsupervised} contains the PCA, MDS/t-SNE and clustering. The statistics
are implemented in \texttt{src/eda.py} and \texttt{src/\allowbreak{}unsupervised.\allowbreak{}py}; this section presents the
graphical results, and their interpretation is in the PRDL report.

\subsection*{Target, age and leakage}

The classes are moderately imbalanced (\NLGG{} LGG, \NGBM{} GBM; majority-class accuracy \MajorityAcc).
Age differs strongly between grades (Figure~\ref{fig:mllb-age}): the empirical cumulative distribution
functions are separated by a Kolmogorov--Smirnov distance of \AgeKS, and the best single threshold
(Youden's index) is \YoudenAge{} years, with sensitivity \YoudenSens{} and specificity \YoudenSpec{} on
the same data (optimistic by construction). Q--Q plots and Shapiro--Wilk tests (in the notebook) show
non-normal age distributions, which is why the Mann--Whitney test is used.

\begin{figure}[H]
\centering
\includegraphics[width=\textwidth]{eda_age_by_grade.png}
\caption{Age at diagnosis by grade: histogram, box and violin plot, and ECDF.}
\label{fig:mllb-age}
\end{figure}

\subsection*{Relevance and redundancy of the predictors}

Figure~\ref{fig:mllb-relevance} ranks every predictor by its AUC used alone, with 2000-resample
stratified bootstrap intervals (\texttt{evaluation.\allowbreak{}bootstrap\_\allowbreak{}ci} with a rank-based AUC written in
NumPy so that it can be resampled quickly). The correlation matrix (Figure~\ref{fig:mllb-corr}) is
drawn with a diverging blue--grey--red scale centred at zero; the maximum VIF is \MaxVIFEDA.
Figure~\ref{fig:mllb-condprob} shows the empirical $P(\mathrm{GBM})$ by age quintile and with and
without each mutation, with Wilson 95\% intervals.

\begin{figure}[H]
\centering
\begin{minipage}[t]{0.49\textwidth}\centering
\includegraphics[width=\textwidth]{eda_feature_relevance.png}
\caption{Single-feature AUC with 95\% bootstrap intervals (colour: direction of the effect).}
\label{fig:mllb-relevance}
\end{minipage}\hfill
\begin{minipage}[t]{0.49\textwidth}\centering
\includegraphics[width=\textwidth]{eda_feature_correlation.png}
\caption{Spearman correlation between predictors (phi coefficient for binary pairs).}
\label{fig:mllb-corr}
\end{minipage}
\end{figure}

\begin{figure}[H]
\centering
\includegraphics[width=\textwidth]{eda_conditional_probability.png}
\caption{Empirical probability of GBM by age quintile and with/without each mutation (95\% Wilson
intervals; dashed line: overall proportion of GBM).}
\label{fig:mllb-condprob}
\end{figure}

\subsection*{Quick baseline and leakage check}

A logistic regression in a \texttt{Pipeline} (scaling and one-hot encoding inside the folds), with the
same stratified 10-fold split used later, checks how much signal the EDA suggests
(Table~\ref{tab:mllb-baseline}). Adding the one-hot diagnosis gives a perfect but fake result.

\begin{table}[H]
\centering
\caption{Quick cross-validated baseline (stratified 10-fold CV, seed 42).}
\label{tab:mllb-baseline}
\resizebox{0.9\textwidth}{!}{\input{tables/eda_baseline.tex}}
\end{table}

\subsection*{Regression target}

Figure~\ref{fig:mllb-agetarget} shows the distribution of age and its relation with the eight
predictors with the strongest Mann--Whitney effects (\NAgeSigBH{} of \NAgeTests{} are significant after
Benjamini--Hochberg correction).

\begin{figure}[H]
\centering
\includegraphics[width=\textwidth]{eda_age_regression_target.png}
\caption{Distribution of age at diagnosis and age with/without the eight most related predictors.}
\label{fig:mllb-agetarget}
\end{figure}

\subsection*{Unsupervised analysis of the mutation profiles}

PCA was computed on centred and on standardised genes (Figure~\ref{fig:mllb-pca}); the biplot arrows
use a single scale factor for both axes so that their directions are not distorted. MDS and t-SNE use
a precomputed Jaccard distance matrix (Figure~\ref{fig:mllb-mds}). The number of clusters was chosen
with the criteria of Figure~\ref{fig:mllb-k}; the rule fixed in advance (highest $k$-means--Ward
agreement) gives $k=\KMain$. Table~\ref{tab:mllb-linkage} and Figure~\ref{fig:mllb-dendro} document why
Jaccard-based hierarchical clustering was not used for the final comparison, and
Figure~\ref{fig:mllb-clugrade} compares the clusters with the grade.

\begin{figure}[H]
\centering
\includegraphics[width=\textwidth]{unsup_pca_biplot.png}
\caption{PCA biplots of the 20 mutation indicators, centred (left) and standardised (right); points
coloured by grade after the PCA.}
\label{fig:mllb-pca}
\end{figure}

\begin{figure}[H]
\centering
\includegraphics[width=\textwidth]{unsup_mds_tsne_grade.png}
\caption{Metric MDS (stress-1 = \MDSStress) and t-SNE maps with Jaccard distance, coloured by grade.}
\label{fig:mllb-mds}
\end{figure}

\begin{figure}[H]
\centering
\includegraphics[width=\textwidth]{unsup_k_selection.png}
\caption{Choice of the number of clusters: within-cluster variation, silhouette, and agreement (ARI)
between $k$-means and Ward.}
\label{fig:mllb-k}
\end{figure}

\begin{table}[H]
\centering
\caption{Hierarchical clustering cut into $k=\KMain$ clusters for each distance and linkage.}
\label{tab:mllb-linkage}
\resizebox{\textwidth}{!}{\input{tables/unsup_linkage_comparison.tex}}
\end{table}

\begin{figure}[H]
\centering
\includegraphics[width=\textwidth]{unsup_dendrograms.png}
\caption{Dendrograms (last 30 merges) with Jaccard distance and average linkage (left) and Euclidean
distance and Ward linkage (right).}
\label{fig:mllb-dendro}
\end{figure}

\begin{figure}[H]
\centering
\includegraphics[width=\textwidth]{unsup_clusters_vs_grade.png}
\caption{Grade within each $k$-means cluster for $k=\KMain$ (main) and $k=\KSecond$ (exploratory).}
\label{fig:mllb-clugrade}
\end{figure}

% =====================================================================
% 4. DEVELOPING MACHINE LEARNING MODELS
% =====================================================================
\section{Developing Machine Learning models}

\subsection{Model Selection}

Each model is defined in \texttt{src/models.py} as a \texttt{Pipeline} with two steps,
\texttt{prep} (a \texttt{Column\allowbreak{}Transformer} from \texttt{src/\allowbreak{}pipelines.\allowbreak{}py}) and \texttt{model}.
Linear models and support vector machines receive standardised inputs and one-hot encoded race
(\texttt{handle\_\allowbreak{}unknown="ignore"}, needed because one race category has a single patient and may be
absent from a training fold); trees and ensembles receive unscaled inputs. Naive Bayes has its own
preprocessing: age is discretised into quantile bins (\texttt{KBinsDiscretizer}) so that a
\texttt{BernoulliNB} can model every column. Tables~\ref{tab:mllb-reg-grids} and~\ref{tab:mllb-clf-grids}
list the models and the grids searched; the grids are small (at most 27 configurations) because
nested cross-validation multiplies the number of fits by 30.

\begin{table}[H]
\centering
\caption{Regression models and hyperparameter grids.}
\label{tab:mllb-reg-grids}
\resizebox{\textwidth}{!}{\input{tables/reg_models_grids.tex}}
\end{table}

\begin{table}[H]
\centering
\caption{Classification models and hyperparameter grids (the SVM grid is a list of two sub-grids,
linear and RBF).}
\label{tab:mllb-clf-grids}
\resizebox{\textwidth}{!}{\input{tables/clf_models_grids.tex}}
\end{table}

\subsection{Model Tuning}

Tuning is done by \texttt{GridSearchCV} with 3 inner folds (stratified for classification,
\texttt{n\_jobs=-1}), scoring \texttt{roc\_auc} for classification and
\texttt{neg\_\allowbreak{}root\_\allowbreak{}mean\_\allowbreak{}squared\_\allowbreak{}error} for regression, inside each of the 10 outer folds
(\texttt{evaluation.\allowbreak{}nested\_\allowbreak{}cv}). Because each outer fold runs its own search, the chosen values can
be compared across folds: a value chosen in most folds indicates a stable choice
(Tables~\ref{tab:mllb-reg-hp} and~\ref{tab:mllb-clf-hp}). The final models used for interpretation are
obtained by \texttt{evaluation.\allowbreak{}fit\_\allowbreak{}final}, which runs the same search once on all patients; its result
is the last column. For logistic regression the final search chose $C=\LRFinalC$ although the folds
chose smaller values; Table~\ref{tab:mllb-flat} shows why: \LRFlatWithin{} of the \LRFlatTotal{}
configurations are within 0.005 AUC of the best one.

\begin{table}[H]
\centering
\caption{Regression: hyperparameters chosen in the 10 outer folds (count in brackets) and in the final
model.}
\label{tab:mllb-reg-hp}
\resizebox{\textwidth}{!}{\input{tables/reg_hyperparameters.tex}}
\end{table}

\begin{table}[H]
\centering
\caption{Classification: hyperparameters chosen in the 10 outer folds and in the final model.}
\label{tab:mllb-clf-hp}
\resizebox{\textwidth}{!}{\input{tables/clf_hyperparameters.tex}}
\end{table}

\begin{table}[H]
\centering
\caption{Flatness of the inner-CV surface of the final classification searches.}
\label{tab:mllb-flat}
\resizebox{\textwidth}{!}{\input{tables/clf_search_flatness.tex}}
\end{table}

A reproducibility problem was found at this stage: two runs of the classification notebook chose
different hyperparameters for logistic regression, because the \texttt{liblinear} solver shuffles the
data and its \texttt{random\_state} had not been fixed. After fixing it, two complete runs give
identical out-of-fold probabilities, and the author's run reproduces every value of the JSON summaries
(Appendix A).

\subsection{Model Evaluation}

\textbf{Regression.} Figure~\ref{fig:mllb-reg-comp} shows the out-of-fold MAE and $R^2$ of every model
with bootstrap intervals (all models between \RegMinMAE{} and \RegMaxMAE{} years of MAE, against
\RegBaseMAE{} for the mean baseline), and Table~\ref{tab:mllb-reg-paired} the paired differences. The
multiple linear regression was refitted with statsmodels for its diagnostics
(Figure~\ref{fig:mllb-ols}); one patient with leverage 1 (the only patient in one race category) is
excluded from the plots because its residual is zero by construction. Figure~\ref{fig:mllb-lasso} shows
the Lasso coefficient path with the chosen penalty ($\alpha=\LassoAlpha$, \LassoNZero{} of
\LassoNCoef{} coefficients set to zero), Figure~\ref{fig:mllb-regtree} the final regression tree
(depth \RegTreeDepth, \RegTreeLeaves{} leaves) and Figure~\ref{fig:mllb-regimp} the permutation
importance computed on the outer test folds.

\begin{figure}[H]
\centering
\includegraphics[width=\textwidth]{reg_model_comparison.png}
\caption{Regression: out-of-fold MAE and $R^2$ with 95\% bootstrap intervals.}
\label{fig:mllb-reg-comp}
\end{figure}

\begin{table}[H]
\centering
\caption{Regression: paired bootstrap of the MAE difference (years) against the mean baseline and
against the model with the lowest MAE.}
\label{tab:mllb-reg-paired}
\small
\input{tables/reg_paired_differences.tex}
\end{table}

\begin{figure}[H]
\centering
\includegraphics[width=0.85\textwidth]{reg_ols_diagnostics.png}
\caption{Residual diagnostics of the multiple linear regression (all patients).}
\label{fig:mllb-ols}
\end{figure}

\begin{figure}[H]
\centering
\begin{minipage}[t]{0.49\textwidth}\centering
\includegraphics[width=\textwidth]{reg_lasso_path.png}
\caption{Lasso path (dotted line: chosen $\alpha$).}
\label{fig:mllb-lasso}
\end{minipage}\hfill
\begin{minipage}[t]{0.49\textwidth}\centering
\includegraphics[width=\textwidth]{reg_permutation_importance.png}
\caption{Regression: permutation importance on the outer test folds.}
\label{fig:mllb-regimp}
\end{minipage}
\end{figure}

\begin{figure}[H]
\centering
\includegraphics[width=\textwidth]{reg_decision_tree.png}
\caption{Final regression tree (node colour: darker = older predicted age).}
\label{fig:mllb-regtree}
\end{figure}

\textbf{Classification.} Figure~\ref{fig:mllb-clf-comp} shows the out-of-fold AUC and balanced
accuracy, Table~\ref{tab:mllb-clf-acc} the accuracy and the per-fold AUC, and
Table~\ref{tab:mllb-clf-paired} the paired AUC differences with the best model. The confusion matrices
and the sensitivity--specificity trade-off as the threshold moves are in Figure~\ref{fig:mllb-confusion};
the threshold is left at 0.5, because tuning it on the same predictions would bias the reported
metrics. Figure~\ref{fig:mllb-clftree} shows the final pruned tree (\texttt{ccp\_alpha} $=\ClfTreeAlpha$,
\ClfTreeCriterion{} criterion), recoloured with the project colours (sklearn's default paints LGG in
orange), and Figures~\ref{fig:mllb-logit}--\ref{fig:mllb-clfimp} the logistic coefficients and the
permutation importance.

\begin{figure}[H]
\centering
\includegraphics[width=\textwidth]{clf_model_comparison.png}
\caption{Classification: out-of-fold AUC and balanced accuracy with 95\% bootstrap intervals.}
\label{fig:mllb-clf-comp}
\end{figure}

\begin{table}[H]
\centering
\begin{minipage}[t]{0.52\textwidth}\centering
\caption{Accuracy (95\% CI) and AUC per outer fold (mean $\pm$ SD).}
\label{tab:mllb-clf-acc}
\resizebox{\textwidth}{!}{\input{tables/clf_accuracy_folds.tex}}
\end{minipage}\hfill
\begin{minipage}[t]{0.45\textwidth}\centering
\caption{Paired bootstrap of the AUC difference with the best model.}
\label{tab:mllb-clf-paired}
\resizebox{\textwidth}{!}{\input{tables/clf_paired_differences.tex}}
\end{minipage}
\end{table}

\begin{figure}[H]
\centering
\includegraphics[width=\textwidth]{clf_confusion_threshold.png}
\caption{Out-of-fold confusion matrices at threshold 0.5 and sensitivity/specificity as a function of
the threshold for the model with the highest AUC.}
\label{fig:mllb-confusion}
\end{figure}

\begin{figure}[H]
\centering
\includegraphics[width=0.9\textwidth]{clf_decision_tree.png}
\caption{Final pruned classification tree (blue = LGG, orange = GBM; darker = purer node).}
\label{fig:mllb-clftree}
\end{figure}

\begin{figure}[H]
\centering
\begin{minipage}[t]{0.49\textwidth}\centering
\includegraphics[width=\textwidth]{clf_logistic_coefficients.png}
\caption{Final logistic regression: coefficients on standardised inputs.}
\label{fig:mllb-logit}
\end{minipage}\hfill
\begin{minipage}[t]{0.49\textwidth}\centering
\includegraphics[width=\textwidth]{clf_permutation_importance.png}
\caption{Classification: permutation importance on the outer test folds.}
\label{fig:mllb-clfimp}
\end{minipage}
\end{figure}

\textbf{Multiclass classification (notebook 06).} The target is built by \texttt{data.\allowbreak{}diagnosis\_\allowbreak{}family},
which maps the six TCGA diagnoses to four families with the dictionary \texttt{config.\allowbreak{}DIAGNOSIS\_\allowbreak{}FAMILY}
and raises an error if a diagnosis is not mapped (Table~\ref{tab:mllb-mc-classes}). The same
\texttt{nested\_cv} function is used with \texttt{task="multiclass"}: the folds are stratified by family,
and the out-of-fold prediction of every patient is a row of four probabilities. The inner criterion is
sklearn's \texttt{roc\_auc\_ovo} (Hand--Till macro AUC); for the 2000 bootstrap resamples the same
metric is recomputed by \texttt{evaluation.\allowbreak{}multiclass\_\allowbreak{}auc}, written with the rank-based binary AUC of the
project and checked to give exactly the value of \texttt{roc\_\allowbreak{}auc\_\allowbreak{}score(multi\_\allowbreak{}class="ovo")}. The
multinomial logistic regression uses the \texttt{saga} solver, the only one in scikit-learn~1.6 that
supports both penalties for the multinomial model. Table~\ref{tab:mllb-mc-grids} lists the models,
Table~\ref{tab:mllb-mc-hp} the chosen hyperparameters and Table~\ref{tab:mllb-mc-paired} the paired
comparison; Figures~\ref{fig:mllb-mc-comp} and~\ref{fig:mllb-mc-imp} show the metrics and the
permutation importance (drop in macro AUC).

\begin{table}[H]
\centering
\caption{Multiclass target: TCGA diagnoses grouped into histological families.}
\label{tab:mllb-mc-classes}
\small
\input{tables/mc_classes.tex}
\end{table}

\begin{table}[H]
\centering
\caption{Multiclass models and hyperparameter grids.}
\label{tab:mllb-mc-grids}
\resizebox{\textwidth}{!}{\input{tables/mc_models_grids.tex}}
\end{table}

\begin{table}[H]
\centering
\caption{Multiclass: hyperparameters chosen in the 10 outer folds and in the final model.}
\label{tab:mllb-mc-hp}
\resizebox{\textwidth}{!}{\input{tables/mc_hyperparameters.tex}}
\end{table}

\begin{figure}[H]
\centering
\includegraphics[width=\textwidth]{mc_model_comparison.png}
\caption{Multiclass: out-of-fold macro AUC and balanced accuracy with 95\% bootstrap intervals.}
\label{fig:mllb-mc-comp}
\end{figure}

\begin{table}[H]
\centering
\caption{Multiclass: paired bootstrap of the macro-AUC difference with the best model.}
\label{tab:mllb-mc-paired}
\small
\input{tables/mc_paired_differences.tex}
\end{table}

\begin{table}[H]
\centering
\caption{Per-family performance of the best multiclass model (out-of-fold; one-vs-rest AUC with 95\%
bootstrap interval).}
\label{tab:mllb-mc-perclass}
\footnotesize
\input{tables/mc_per_class.tex}
\end{table}

\begin{figure}[H]
\centering
\includegraphics[width=0.75\textwidth]{mc_permutation_importance.png}
\caption{Multiclass: permutation importance on the outer test folds.}
\label{fig:mllb-mc-imp}
\end{figure}

\textbf{SHAP values (notebook 07).} The notebook reads the JSON summaries of notebooks 05 and 06 to find
the winning models (the best one and those whose paired difference with it includes zero), refits them
with \texttt{fit\_final} and asserts that the hyperparameters are the ones saved by those notebooks.
\texttt{explain.\allowbreak{}shap\_\allowbreak{}values} applies the pipeline's preprocessing and chooses the explainer by model
type (\texttt{TreeExplainer} for forests, trees and XGBoost; \texttt{LinearExplainer} with all patients
as background for logistic regression); it returns the values, the base value, the units and the raw
model output, and \texttt{explain.\allowbreak{}additivity\_\allowbreak{}error} checks that base value plus contributions reproduces
every prediction. The summary plots (\texttt{explain.\allowbreak{}beeswarm}) are drawn with Matplotlib rather than
with the library's own plots, to keep the project style. Table~\ref{tab:mllb-shap-agree} compares the
SHAP ranking with the permutation importance, and Figure~\ref{fig:mllb-shap-age} shows the dependence of
the age contribution.

\begin{table}[H]
\centering
\caption{Agreement between SHAP (share of mean $|$SHAP$|$, race columns summed) and permutation importance.}
\label{tab:mllb-shap-agree}
\resizebox{\textwidth}{!}{\input{tables/shap_binary_agreement.tex}}
\end{table}

\begin{figure}[H]
\centering
\includegraphics[width=\textwidth]{shap_binary_age.png}
\caption{SHAP dependence of age for the three tied binary models (colour: IDH1 status).}
\label{fig:mllb-shap-age}
\end{figure}

\textbf{Causal analysis (notebook 08).} \texttt{causal.\allowbreak{}mediation\_\allowbreak{}decomposition} fits the three least-squares
regressions with \texttt{numpy.\allowbreak{}linalg.\allowbreak{}lstsq} (fast enough for 2000 bootstrap resamples) and returns the
identity check $c-(c'+ab)$, which the notebook asserts to be below $10^{-8}$;
\texttt{causal.\allowbreak{}ols\_\allowbreak{}table} refits the models with statsmodels and HC3 standard errors for the reported
intervals, because the residuals of age are heteroscedastic (Figure~\ref{fig:mllb-ols}). Race enters as a single
white/non-white indicator: with the four categories, the patient alone in one category has leverage 1,
and the HC3 correction divides by $1-h_{ii}$. Table~\ref{tab:mllb-causal-adj} lists the four adjustment
sets and Table~\ref{tab:mllb-causal-int} the interaction model.

\begin{table}[H]
\centering
\caption{Coefficient of IDH mutation on age at diagnosis (years) for different adjustment sets (HC3
intervals).}
\label{tab:mllb-causal-adj}
\small
\input{tables/causal_adjustment_sets.tex}
\end{table}

\begin{table}[H]
\centering
\caption{Age at diagnosis with an IDH $\times$ grade interaction (OLS, HC3 intervals).}
\label{tab:mllb-causal-int}
\small
\input{tables/causal_interaction.tex}
\end{table}

% =====================================================================
% 5. CODE DISTRIBUTION
% =====================================================================
\section{Code Distribution}\label{sec:mllb-code}

All code is in the public repository \url{https://github.com/orianafe16/ML}. The raw data are included
(CC BY 4.0); the course templates and the literature PDFs are excluded with \texttt{.gitignore}.

\begin{lstlisting}[language={}]
ML/
  data/raw/          original CSV files (unchanged)
  data/processed/    clean dataset written by notebook 01 (not versioned)
  src/               config, data, eda, unsupervised, pipelines, models,
                     evaluation, explain, causal, plotting, reporting,
                     check_env
  notebooks/         01_data_preparation ... 08_causal
  results/           figures/, tables/ (LaTeX), logs/ (JSON per stage,
                     out-of-fold predictions)
  report/            PRDL and MLLB sources, references.bib, make_numbers.py,
                     build_zips.py, check_references.py
  docs/              DECISIONS.md (justification of every decision),
                     AI_USAGE_LOG.md
  environment.yml, requirements.txt, README.md
\end{lstlisting}

To reproduce every result:
\begin{lstlisting}[language={}]
conda env create -f environment.yml
conda activate glioma-ml
python -m ipykernel install --user --name glioma-ml
python -m src.check_env
# run notebooks/01 ... 08 in order with the "glioma-ml" kernel
python report/make_numbers.py
python report/build_zips.py       # Overleaf bundles of both reports
\end{lstlisting}

\texttt{docs/\allowbreak{}DECISIONS.\allowbreak{}md} records every design decision with its justification and the course material
it is based on, and \texttt{docs/\allowbreak{}AI\_\allowbreak{}USAGE\_\allowbreak{}LOG.\allowbreak{}md} every interaction with AI assistants.

% =====================================================================
% 6. CONCLUSIONS
% =====================================================================
\section{Conclusions}

From the point of view of tools and methodology:
\begin{itemize}
  \item \textbf{The pipeline is reproducible end to end.} The rebuilt dataset matches the UCI file in
        every value; two complete runs give identical results; and the run on a second computer (the
        author's) reproduced every value of the JSON summaries. The only reproducibility failure
        found (an unseeded solver) was detected precisely by this comparison.
  \item \textbf{Nested cross-validation is affordable at this scale.} With small grids, each notebook
        runs in a few minutes, and keeping the chosen hyperparameters of every outer fold gives a direct
        view of the stability of the tuning.
  \item \textbf{Results.} In classification, all models reach out-of-fold AUCs between \ClfMinAUC{} and
        \ClfMaxAUC{} (baseline \ClfAUCBase); in regression, MAEs between \RegMinMAE{} and \RegMaxMAE{}
        years (baseline \RegBaseMAE). Under the comparison rule fixed in advance no model is declared
        better than the others in either task; the paired comparison separates only small differences
        in classification. The multiclass task reaches a macro AUC of \McBestAUC{} for the \McBest{}
        (baseline \McBaseAUC). The interpretation of these results is in the PRDL report.
  \item \textbf{Reusing the same machinery.} The multiclass task needed only a new metric function and
        a new branch in \texttt{nested\_cv}; the SHAP notebook re-fits the final models with
        \texttt{fit\_final} and asserts that it obtains the hyperparameters saved by notebooks 05 and 06,
        and every SHAP explanation is checked for additivity before it is plotted.
  \item \textbf{Graphics that check themselves.} Fixing one colour per class across all figures, and
        inspecting every figure before accepting it, caught errors that no test would have detected: a
        biplot with distorted arrows, a bar chart whose axis hid part of the data, a tree with inverted
        class colours and a tree rule printed as ``IDH1 $\leq$ 0.0''.
  \item \textbf{Generating the reports from the results} (numbers as macros, tables as \texttt{.tex}
        files, figures copied by a script) means that re-running the notebooks updates both reports
        without manual editing.
\end{itemize}



% =====================================================================
% REFERENCES
% =====================================================================
\section{References}
\begingroup
% thebibliography opens its own \section*{References}; it is swallowed here so that
% the numbered banner above is the only heading.
\renewcommand{\section}[2]{}
\bibliographystyle{IEEEtran}
\bibliography{references}
\endgroup

\newpage

% =====================================================================
% APPENDIX: ENSEMBLE LEARNING EXPERIMENT
% =====================================================================
\appendixbanner{APPENDIX A: ENSEMBLE LEARNING EXPERIMENT}

\subsection*{What we did}

As the applied part of our Ensemble Learning topic, we compared a single Decision Tree
against four ensemble methods (Random Forest, AdaBoost, Gradient Boosting, and a Stacking
ensemble) on the Breast Cancer Wisconsin dataset (569 patients, 30 features), framed as a
clinical screening problem where missing a malignant case matters more than a false
alarm. All five models were trained on a stratified 75/25 split (426 train / 143 test,
\texttt{random\_state=42}) and evaluated with accuracy, sensitivity, specificity,
F1-score, balanced accuracy, and ROC-AUC, ranking models by ROC-AUC and balanced accuracy
together rather than by accuracy alone. We then tuned the decision threshold of the
top-ranked model to see what it takes to reach full sensitivity.

\subsection*{Results}

All four ensembles clearly beat the single Decision Tree (ROC-AUC 0.938), with Gradient
Boosting and Random Forest essentially tied for the best ROC-AUC (0.9950 and 0.9949) and
AdaBoost reaching the best balanced accuracy (0.962) and the only perfect specificity
(1.000) of the five models.

\begin{table}[h]
\centering
\caption{Model comparison at the default 0.5 threshold}
\label{tab:default-results}
\begin{tabularx}{\textwidth}{l r r r r}
\toprule
\textbf{Model} & \textbf{Sensitivity} & \textbf{Specificity} & \textbf{Balanced acc.} & \textbf{ROC-AUC} \\
\midrule
Decision Tree & 0.9245 & 0.9556 & 0.9400 & 0.9384 \\
Random Forest & 0.9245 & 0.9778 & 0.9512 & 0.9949 \\
AdaBoost & 0.9245 & 1.0000 & 0.9623 & 0.9897 \\
Gradient Boosting & 0.9057 & 0.9889 & 0.9473 & 0.9950 \\
Stacking & 0.9245 & 0.9778 & 0.9512 & 0.9929 \\
\bottomrule
\end{tabularx}
\end{table}

The more useful finding was methodological rather than numerical. Gradient boosting,
the top-ranked model by ROC-AUC, needed an aggressive threshold of 0.005 to catch every
malignant case in the test set (10 false positives out of 90 benign cases), while random
forest, whose ROC-AUC is only 0.0001 lower, reached the same full sensitivity at a much
gentler threshold of 0.3.

\begin{figure}[h]
\centering
\includegraphics[width=0.7\textwidth]{figures/roc_curves.png}
\caption{ROC curves for all five models on the held-out test set}
\label{fig:roc-curves}
\end{figure}

\begin{figure}[h]
\centering
\includegraphics[width=0.6\textwidth]{figures/threshold_tradeoff.png}
\caption{Sensitivity/specificity trade-off as the decision threshold decreases (Gradient Boosting)}
\label{fig:threshold-tradeoff}
\end{figure}

\subsection*{Conclusion}

A higher ROC-AUC does not by itself pick the model best suited to a screening task: the
top-ranked model by that single metric was also the hardest one to operate safely,
requiring a far more extreme threshold than a near-tied alternative to reach the same
sensitivity. Ranking by balanced accuracy alongside ROC-AUC, and checking the threshold
each model actually needs, gives a more honest picture than accuracy or ROC-AUC alone.

\subsection*{Code}

The full, reproducible notebook (\texttt{ensemble\_learning\_experiment.ipynb}) is
available in our group's repository: \texttt{https://github.com/orianafe16/PRDL-MLLB.git}.

\newpage
% =====================================================================
% APPENDIX B: USE OF AI CODE ASSISTANTS
% =====================================================================
\appendixbanner{APPENDIX B: USE OF AI CODE ASSISTANTS}
\addcontentsline{toc}{section}{Appendix B: Use of AI code assistants}

\subsection*{Workflow}

The code was co-developed with Claude (Anthropic) working on the project folder through the Claude
desktop app. The division of work was fixed from the start: the author defined the dataset, the tasks,
the models and the validation protocol (listed in \texttt{docs/DECISIONS.md} with their justification);
the assistant wrote the code stage by stage, executed every notebook in a clean copy of the repository
and inspected every figure before delivering it; the author ran each notebook on her own computer, and
the assistant compared the resulting JSON summaries with its own run, value by value, before the next
stage. Every exchange that shaped the code is recorded in \texttt{docs/AI\_USAGE\_LOG.md}.

\subsection*{Errors detected (candidates for the AI Error Bounty)}

All of them were produced by the assistant and caught by the verification loop before delivery.

\begin{enumerate}
  \item \textbf{Non-reproducible results.} \texttt{LogisticRegression(solver="liblinear")} was created
        without \texttt{random\_state}; \texttt{liblinear} shuffles the data, so two runs chose different
        hyperparameters in the outer folds. \emph{Detection}: comparing two runs. \emph{Fix}:
        \texttt{random\_state=42}; two complete runs then gave identical out-of-fold probabilities.
  \item \textbf{Silent NaN in the regression diagnostics.} The summary reported ``maximum Cook's
        distance = nan''. \emph{Cause}: one race category has a single patient, so its one-hot column
        gives that patient leverage 1, the model fits it exactly, and the standardised residual is 0/0.
        \emph{Fix}: detect leverage-1 patients, report them and exclude them from the diagnostics.
  \item \textbf{Distorted biplot.} The loading arrows were scaled with a different factor on each axis,
        which changes their directions. \emph{Detection}: arrows missing from the scaled biplot.
        \emph{Fix}: one common factor.
  \item \textbf{Misleading tree figures.} sklearn's default colours painted LGG in orange and GBM in
        blue, the opposite of every other figure; and with zero decimals the split ``IDH1 $\leq$ 0.5''
        was printed as ``IDH1 $\leq$ 0.0''. \emph{Detection}: comparing the figure with the text export of
        the tree. \emph{Fix}: recolouring and one decimal.
  \item \textbf{Axis hiding data.} The number of mutations per patient was drawn on a categorical axis
        that silently skipped empty categories, hiding the few patients with many mutations.
        \emph{Fix}: grouped ``6+'' category.
  \item \textbf{Escaping errors in generated tables} (\texttt{\$\textbackslash pm\$} printed literally)
        and minor bugs (a negative error-bar length from floating-point rounding, a lost column name,
        NumPy~2 number formatting) found when executing the notebooks.
\end{enumerate}

\subsection*{Reflection}

The assistant wrote code that ran without errors and looked plausible in every one of these cases; the
errors were only visible by running the code twice, reading the figures against the numbers, or
connecting a NaN with the theory of high-leverage points. Owning the code meant owning this
verification loop.


\end{document}
